\chapter{Kiến trúc Hệ thống và Thuật toán EAGNG Phân cụm Khuyết tật Trực tuyến}
\label{chap:c3}

Chương \ref{chap:c3} trình bày chi tiết về kiến trúc hệ thống tổng thể và thiết kế thuật toán Mạng khí Nơ-ron tăng trưởng tích lũy bằng chứng đề xuất cho bài toán tự động phân loại khuyết tật màn hình điện thoại. Giải pháp được xây dựng dựa trên mô hình hai pha nối tiếp nhằm giải quyết đồng thời yêu cầu lọc bỏ sản phẩm đạt chuẩn ở tốc độ cao và phân cụm linh hoạt các dạng lỗi bề mặt phức tạp trong điều kiện vận hành thời gian thực. Hệ thống tích hợp thông tin entropy cục bộ làm kênh dữ liệu thứ tư kết hợp với các mạng nơ-ron sâu chuyên biệt để trích xuất đặc trưng phù hợp cho từng giai đoạn xử lý.

Nội dung chính của chương tập trung làm rõ cấu trúc toán học và các cơ chế cải tiến cốt lõi của thuật toán đề xuất. Pha 1 sử dụng mạng mã hóa tự động để nén không gian biểu diễn và thiết lập ranh giới phân loại bất thường bằng quy tắc ba lần độ lệch chuẩn trên Mạng khí Nơ-ron tăng trưởng ngoại tuyến. Pha 2 phát triển thuật toán Mạng khí Nơ-ron tăng trưởng tích lũy bằng chứng với bốn cơ chế cải tiến chính bao gồm: (i) suy giảm tốc độ học và tích lũy lỗi theo tuổi nút, (ii) trì hoãn chèn nút dựa trên tích lũy bằng chứng kết hợp huấn luyện cục bộ, (iii) phát triển cấu trúc nhạy mật độ, và (iv) kiểm tra sự chồng lấp hộp giới hạn để kiểm soát liên kết cạnh. Toàn bộ quy trình tạo thành một khung xử lý toán học hoàn chỉnh, làm cơ sở lý thuyết cho việc triển khai thực nghiệm ở chương tiếp theo.

\section{Kiến trúc tổng thể hệ thống đề xuất}
\label{sec:kientruc_tongthe}

Hệ thống đề xuất được cấu trúc thành hai pha xử lý nối tiếp nhằm giải quyết triệt để bài toán phát hiện bất thường và phân cụm khuyết tật thời gian thực trên màn hình điện thoại. Pha 1 đảm nhận vai trò lọc bỏ các mẫu sản phẩm đạt chuẩn thông qua mô hình phát hiện bất thường ngoại tuyến, giúp giảm tải khối lượng tính toán cho hệ thống. Pha 2 tập trung phân tích hình thái khuyết tật và thực hiện phân cụm tự động bằng thuật toán Mạng khí Nơ-ron tăng trưởng tích lũy bằng chứng. Sự kết hợp giữa xử lý dữ liệu đa kênh và trích xuất đặc trưng chuyên biệt giữa hai pha tạo nên một quy trình khép kín, đảm bảo độ chính xác và tính thích nghi cao trong môi trường sản xuất công nghiệp. Sơ đồ khối tổng thể của hệ thống được trình bày chi tiết trong Hình \ref{fig:kientruc_tongthe}.

\begin{figure}[H]
\centering
\begin{tikzpicture}[
    node distance=1.2cm and 0.5cm,
    block/.style={draw=black!80, fill=blue!5, rectangle, rounded corners=2pt, minimum width=1.7cm, minimum height=1.0cm, align=center, font=\normalsize},
    phase/.style={draw=black!60, fill=gray!10, rectangle, dashed, inner sep=6pt},
    arrow/.style={->, >=stealth, thick, color=black!80},
    font=\small
]

% Phase 1 Nodes
\node [block] (input1) {Ảnh RGB\\đạt chuẩn};
\node [block, right=of input1] (entropy1) {Tính toán\\Entropy};
\node [block, right=of entropy1] (ae) {Mạng mã hóa\\tự động};
\node [block, right=of ae] (gng_off) {Mạng GNG\\ngoại tuyến};
\node [block, right=of gng_off] (check3s) {Kiểm tra\\ngưỡng $3\sigma$};

% Phase 2 Nodes
\node [block, below=1.8cm of check3s] (resnet) {Mạng ResNet\\bốn kênh};
\node [block, left=of resnet] (eagng) {Thuật toán\\EAGNG};
\node [block, left=of eagng] (cc) {Thành phần\\liên kết};
\node [block, left=of cc] (output) {Phân loại\\cụm lỗi};

% Arrows Phase 1
\draw [arrow] (input1) -- (entropy1);
\draw [arrow] (entropy1) -- (ae);
\draw [arrow] (ae) -- (gng_off);
\draw [arrow] (gng_off) -- (check3s);

% Decision & Transition
\draw [arrow] (check3s) -- node[left, font=\small] {Bất thường ($>3\sigma$)} (resnet);

% Arrows Phase 2
\draw [arrow] (resnet) -- (eagng);
\draw [arrow] (eagng) -- (cc);
\draw [arrow] (cc) -- (output);

% Background Phases
\begin{scope}[on background layer]
    \node [phase, fit=(input1) (check3s), label={above:\small Pha 1: Phát hiện bất thường ngoại tuyến}] (p1) {};
    \node [phase, fit=(resnet) (output), label={below:\small Pha 2: Phân cụm lỗi thích nghi trực tuyến bằng thuật toán EAGNG}] (p2) {};
\end{scope}

\end{tikzpicture}
\caption{Sơ đồ kiến trúc tổng thể hai pha của hệ thống tự động phân loại khuyết tật màn hình điện thoại.}
\label{fig:kientruc_tongthe}
\end{figure}

\subsection{Pha 1: Phát hiện bất thường ngoại tuyến}
\label{subsec:pha1_overview}

Pha phát hiện bất thường ngoại tuyến có mục tiêu thiết lập mô hình biểu diễn phân bố của sản phẩm đạt tiêu chuẩn chất lượng. Quy trình khởi đầu bằng việc thu thập tập hợp các hình ảnh màn hình điện thoại không chứa khuyết tật từ dây chuyền sản xuất. Mỗi hình ảnh đầu vào ban đầu được tính toán bản đồ entropy cục bộ để ghi nhận thông tin về độ phức tạp bề mặt và được ghép nối trực tiếp thành kênh dữ liệu thứ tư. Mạng mã hóa tự động tiếp nhận dữ liệu bốn kênh này nhằm thực hiện nén dữ liệu và thu gọn không gian biểu diễn về một không gian tiềm ẩn có kích thước nhỏ hơn \cite{hinton2006reducing}.

Mô hình Mạng khí Nơ-ron tăng trưởng ngoại tuyến được huấn luyện trực tiếp trên tập hợp các vectơ tiềm ẩn thu được từ mạng mã hóa tự động. Quá trình huấn luyện này giúp Mạng khí Nơ-ron tăng trưởng học được cấu trúc không gian của dữ liệu chuẩn và hình thành bộ nút đại diện cho phân bố dữ liệu không lỗi. Khi một hình ảnh mới được đưa vào kiểm tra, vectơ tiềm ẩn của hình ảnh đó sẽ được tính toán khoảng cách Euclid tới nút gần nhất trong Mạng khí Nơ-ron tăng trưởng.

Việc quyết định một mẫu dữ liệu có chứa khuyết tật hay không dựa trên quy tắc thống kê ba lần độ lệch chuẩn. Khoảng cách Euclid cực tiểu $d_{min}(z)$ từ vectơ tiềm ẩn $z$ tới tập hợp các nút trong Mạng khí Nơ-ron tăng trưởng được xác định theo Công thức \ref{eq:distance_p1}:
\begin{equation}
\label{eq:distance_p1}
d_{min}(z) = \min_{i \in V} \| z - w_i \|_2
\end{equation}
trong đó $z$ đại diện cho vectơ tiềm ẩn thu được từ không gian ẩn của mạng mã hóa tự động, $V$ biểu diễn tập hợp các nút trong đồ thị nơ-ron khí tăng trưởng ngoại tuyến, $w_i$ là vectơ trọng số của nút thứ $i$ trong tập hợp $V$, và $\| \cdot \|_2$ thể hiện chuẩn Euclid hai. Ngưỡng quyết định bất thường $T_{pha1}$ được thiết lập theo Công thức \ref{eq:threshold_p1}:
\begin{equation}
\label{eq:threshold_p1}
T_{pha1} = \mu_d + 3 \sigma_d
\end{equation}
trong đó $\mu_d$ biểu diễn giá trị trung bình của khoảng cách cực tiểu trên tập dữ liệu huấn luyện chuẩn, và $\sigma_d$ đại diện cho độ lệch chuẩn tương ứng của chuỗi khoảng cách này. Nếu khoảng cách $d_{min}(z)$ vượt quá ngưỡng $T_{pha1}$, mẫu dữ liệu được xác định là bất thường và được chuyển sang pha 2 để thực hiện phân cụm khuyết tật.

\subsection{Pha 2: Phân cụm lỗi thích nghi trực tuyến bằng thuật toán EAGNG}
\label{subsec:pha2_overview}

Pha phân cụm lỗi thích nghi trực tuyến tiếp nhận toàn bộ các mẫu hình ảnh được xác định là bất thường từ pha 1 để tiến hành phân tích sâu và phân loại cấu trúc lỗi. Không giống như pha 1 tập trung vào việc tái tạo toàn cục thông qua mạng mã hóa tự động, pha 2 sử dụng kiến trúc mạng nơ-ron tích chập ResNet điều chỉnh bốn kênh đầu vào để trích xuất các đặc trưng chú trọng đến hình thái và ranh giới khuyết tật.

Vectơ đặc trưng thu được từ mạng ResNet bốn kênh được đưa trực tiếp vào mô hình Mạng khí Nơ-ron tăng trưởng tích lũy bằng chứng để tiến hành huấn luyện trực tuyến. Mạng khí Nơ-ron tăng trưởng tích lũy bằng chứng thực hiện cập nhật cấu trúc mạng một cách liên tục dựa trên bốn cơ chế cải tiến chính bao gồm: (i) suy giảm tốc độ học và tích lũy lỗi theo tuổi nút, (ii) trì hoãn chèn nút dựa trên tích lũy bằng chứng kết hợp huấn luyện cục bộ, (iii) phát triển cấu trúc nhạy mật độ nhằm kiểm soát sự bùng nổ số lượng nút, và (iv) kiểm tra sự chồng lấp hộp giới hạn để kiểm soát việc nối cạnh giữa các cụm. Các cơ chế này giúp mô hình duy trì tính ổn định cấu trúc tô-pô và tránh hiện tượng nhiễu cấu trúc khi dữ liệu xuất hiện theo dòng liên tục \cite{fritzke1995growing}.

Sau khi Mạng khí Nơ-ron tăng trưởng tích lũy bằng chứng hoàn thành cập nhật trạng thái trực tuyến cho mẫu dữ liệu, thuật toán thành phần liên kết trên đồ thị được áp dụng để nhóm các nút kết nối thành các cụm khuyết tật riêng biệt. Các tập hợp nút thuộc cùng một thành phần liên kết trên đồ thị được xác định là một lớp khuyết tật cụ thể trên bề mặt màn hình điện thoại. Phương pháp phân cụm dựa trên đồ thị này giúp hệ thống tự động phát hiện các dạng lỗi mới phát sinh mà không cần biết trước số lượng cụm khuyết tật trong quá trình vận hành v.v...
\section{Biểu diễn dữ liệu và Xử lý đặc trưng không gian Entropy}
\label{sec:bieudien_dulieu}

Biểu diễn dữ liệu đa kênh đóng vai trò quyết định trong việc nâng cao độ nhạy của hệ thống đối với các khuyết tật có độ tương phản thấp trên bề mặt màn hình điện thoại. Các phương pháp xử lý ảnh truyền thống dựa trên không gian màu RGB thường gặp khó khăn trong việc phát hiện các vết ố mờ do sự biến đổi độ sáng không đáng kể so với nền xung quanh. Để khắc phục hạn chế này, hệ thống đề xuất tích hợp không gian entropy cục bộ như một kênh thông tin bổ sung nhằm gia tăng khả năng phân tách đặc trưng giữa vùng bề mặt chuẩn và vùng chứa khuyết tật.

\subsection{Thuật toán trích xuất bản đồ Entropy cục bộ}
\label{subsec:entropy_extraction}

Entropy cục bộ đo lường mức độ không chắc chắn hoặc độ phức tạp của sự phân bố mức xám trong một cửa sổ trượt xung quanh mỗi điểm ảnh. Bản đồ entropy thu được phản ánh sự thay đổi cấu trúc bề mặt, trong đó các vùng chứa đốm tròn hoặc hạt bụi thường thể hiện độ mất trật tự thông tin cao hơn đáng kể so với các vùng bề mặt đồng nhất \cite{shannon1948mathematical}.

Quy trình trích xuất bản đồ entropy cục bộ được thực hiện thông qua các bước tính toán trên từng vùng ảnh lân cận. Đầu tiên, hình ảnh đầu vào được chuyển đổi về dạng ảnh mức xám hai chiều. Tiếp theo, một cửa sổ trượt vuông kích thước $N \times N$ điểm ảnh được quét qua toàn bộ bề mặt hình ảnh với bước nhảy bằng một. Tại mỗi vị trí tâm $(x,y)$ của cửa sổ trượt $W(x,y)$, xác suất xuất hiện $p_k$ của mức xám $k$ được xác định theo Công thức \ref{eq:probability}:
\begin{equation}
\label{eq:probability}
p_k = \frac{n_k}{\sum_{j=0}^{L-1} n_j}
\end{equation}
trong đó $n_k$ đại diện cho số lượng điểm ảnh có mức xám $k$ nằm bên trong cửa sổ trượt $W(x,y)$, $L$ biểu diễn tổng số mức xám của hình ảnh với giá trị mặc định bằng 256, và mẫu số thể hiện tổng số điểm ảnh trong cửa sổ trượt $N \times N$.

Giá trị entropy cục bộ $E(x,y)$ tại điểm ảnh tâm $(x,y)$ được tính toán dựa trên định lý entropy thông tin Shannon theo Công thức \ref{eq:entropy}:
\begin{equation}
\label{eq:entropy}
E(x,y) = -\sum_{k=0}^{L-1} p_k \log_2 (p_k + \epsilon)
\end{equation}
trong đó $p_k$ là xác suất xuất hiện của mức xám $k$, $\epsilon$ là một hằng số dương cực tiểu được thêm vào nhằm tránh lỗi vô cấm toán học khi $p_k = 0$, và $\log_2$ là hàm logarit cơ số hai. Việc lặp lại phép tính trên toàn bộ các điểm ảnh giúp tạo ra một bản đồ entropy hai chiều có cùng kích thước không gian với hình ảnh ban đầu v.v...

Việc lựa chọn kích thước cửa sổ trượt $N \times N$ ảnh hưởng trực tiếp đến chất lượng trích xuất đặc trưng hình thái. Nếu kích thước $N$ quá nhỏ, bản đồ entropy sẽ bị ảnh hưởng mạnh bởi nhiễu ngẫu nhiên của cảm biến thu nhận ảnh. Ngược lại, nếu kích thước $N$ quá lớn, ranh giới của các vết khuyết tật siêu nhỏ sẽ bị làm mịn quá mức, dẫn đến mất mát thông tin không gian quan trọng. Qua các thực nghiệm khảo sát, kích thước cửa sổ $N = 9$ điểm ảnh được lựa chọn nhằm tối ưu hóa sự cân bằng giữa khả năng triệt tiêu nhiễu và độ sắc nét của biên giới khuyết tật.

\subsection{Kỹ thuật ghép kênh RGB-E và Chuẩn hóa dữ liệu}
\label{subsec:channel_concatenation}

Kỹ thuật ghép kênh kết hợp bản đồ entropy vào không gian màu nguyên bản nhằm tạo ra cấu trúc dữ liệu bốn kênh RGB-E đồng bộ về kích thước không gian. Cấu trúc dữ liệu mở rộng này cung cấp đầy đủ thông tin màu sắc, độ sáng và cấu trúc bề mặt cho các mô hình trích xuất đặc trưng phía sau.

Vectơ biểu diễn dữ liệu bốn kênh $I_{RGB-E}(x,y)$ tại tọa độ $(x,y)$ được hình thành bằng cách ghép nối trực tiếp bản đồ entropy vào tập hợp các kênh màu tiêu chuẩn theo Công thức \ref{eq:rgbe}:
\begin{equation}
\label{eq:rgbe}
I_{RGB-E}(x,y) = \left[ R(x,y), G(x,y), B(x,y), E_{norm}(x,y) \right]^T
\end{equation}
trong đó $R(x,y)$, $G(x,y)$, và $B(x,y)$ lần lượt đại diện cho cường độ giá trị của các kênh màu đỏ, xanh lá cây và xanh dương tại điểm ảnh $(x,y)$, trong khi $E_{norm}(x,y)$ biểu diễn giá trị entropy cục bộ đã qua chuẩn hóa.

Do phạm vi giá trị của bản đồ entropy ban đầu khác biệt so với khoảng giá trị màu sắc, việc chuẩn hóa dữ liệu là bắt buộc nhằm đảm bảo tính đồng đều về mặt miền giá trị giữa các kênh đầu vào. Giá trị entropy cục bộ $E(x,y)$ được biến đổi về khoảng $[0, 1]$ thông qua kỹ thuật chuẩn hóa Cực tiểu - Cực đại theo Công thức \ref{eq:minmax}:
\begin{equation}
\label{eq:minmax}
E_{norm}(x,y) = \frac{E(x,y) - E_{min}}{E_{max} - E_{min} + \delta}
\end{equation}
trong đó $E_{min}$ và $E_{max}$ lần lượt đại diện cho giá trị entropy nhỏ nhất và lớn nhất ghi nhận được trên toàn bộ bản đồ ảnh, còn $\delta$ là hằng số dương cực tiểu ngăn ngừa hiện tượng chia cho số không.

Quy trình chuẩn hóa bốn kênh giúp ổn định quá trình lan truyền ngược và tăng tốc độ hội tụ khi huấn luyện các mạng nơ-ron sâu trong hệ thống. Dữ liệu bốn kênh sau khi xử lý đáp ứng đầy đủ điều kiện về mặt phân bố để làm đầu vào cho mô hình Autoencoder ở Pha 1 và mạng ResNet cải tiến ở Pha 2 \cite{he2016deep}.
\section{Trích xuất đặc trưng chuyên biệt cho hai pha}
\label{sec:trichxuat_dactrung}

Quá trình trích xuất đặc trưng được thiết kế tách biệt cho hai pha xử lý nhằm đáp ứng các mục tiêu bài toán khác nhau. Pha phát hiện bất thường ngoại tuyến yêu cầu một không gian biểu diễn thu nhỏ để nén các đặc tính của bề mặt màn hình đạt chuẩn, trong khi pha phân cụm lỗi trực tuyến cần các đặc trưng phản ánh chi tiết hình thái và biên giới của khuyết tật. Sự phân tách này giúp tối ưu hóa dung lượng tính toán và nâng cao khả năng phân tách giữa các lớp lỗi phát sinh.

\subsection{Mạng mã hóa tự động: Nén không gian tiềm ẩn nhận dạng mẫu chuẩn}
\label{subsec:autoencoder}

Mạng mã hóa tự động đảm nhận nhiệm vụ học cách biểu diễn thu gọn của các hình ảnh màn hình không chứa khuyết tật. Bằng cách huấn luyện mạng tái tạo lại chính dữ liệu đầu vào bốn kênh, mô hình buộc phải nén thông tin qua một bộ thắt nút để tạo ra vectơ tiềm ẩn chứa các đặc trưng tổng quát nhất của bề mặt chuẩn \cite{kingma2013autoencoding}.

Quá trình mã hóa và giải mã được thực hiện qua hai phép biến đổi phi tuyến nối tiếp. Hàm mã hóa chuyển đổi dữ liệu đầu vào bốn kênh thành vectơ tiềm ẩn theo Công thức \ref{eq:encoder}:
\begin{equation}
\label{eq:encoder}
z = g_\theta(I_{RGB-E})
\end{equation}
trong đó $I_{RGB-E}$ đại diện cho dữ liệu ảnh đầu vào bốn kênh bao gồm ba kênh màu và một kênh entropy, $g_\theta$ biểu diễn hàm mã hóa được tham số hóa bởi tập trọng số $\theta$, và $z$ là vectơ tiềm ẩn nằm trong không gian nén. Sau đó, quá trình tái tạo hình ảnh từ không gian tiềm ẩn được thực hiện qua hàm giải mã theo Công thức \ref{eq:decoder}:
\begin{equation}
\label{eq:decoder}
\hat{I}_{RGB-E} = f_\phi(z)
\end{equation}
trong đó $f_\phi$ là hàm giải mã được tham số hóa bởi tập trọng số $\phi$, và $\hat{I}_{RGB-E}$ biểu diễn dữ liệu ảnh tái tạo bốn kênh thu được ở đầu ra.

Vectơ tiềm ẩn $z$ thu được sau quá trình huấn luyện giữ vai trò là đầu vào cho mô hình Mạng khí Nơ-ron tăng trưởng ngoại tuyến ở Pha 1. Do mạng mã hóa tự động chỉ được huấn luyện trên tập dữ liệu đạt chuẩn, khả năng tái tạo của mạng đối với các vùng chứa khuyết tật sẽ kém hiệu quả, tạo ra sự sai lệch khoảng cách đáng kể và hỗ trợ cho việc phân loại bất thường ở các giai đoạn sau.

\subsection{Mạng ResNet bốn kênh: Trích xuất đặc trưng hình thái lỗi}
\label{subsec:resnet4khenh}

Mạng ResNet bốn kênh được cải tiến để trích xuất các đặc trưng ngữ nghĩa và hình thái khuyết tật phục vụ cho giai đoạn phân cụm trực tuyến ở Pha 2. Khác với mạng mã hóa tự động tập trung vào việc tái tạo toàn cục, kiến trúc ResNet tận dụng các liên kết tắt nhằm bảo toàn thông tin độ dốc và trích xuất các đặc trưng ranh giới sắc nét của đốm tròn hoặc hạt bụi \cite{he2016deep}.

Lớp tích phân đầu tiên của kiến trúc mạng được điều chỉnh để tiếp nhận trực tiếp dữ liệu đầu vào bốn kênh. Biến đổi đặc trưng từ ảnh bất thường sang vectơ đặc trưng hình thái được mô tả theo Công thức \ref{eq:resnet_extract}:
\begin{equation}
\label{eq:resnet_extract}
f_{defect} = R_\psi(I_{RGB-E}^{anomaly})
\end{equation}
trong đó $I_{RGB-E}^{anomaly}$ đại diện cho dữ liệu ảnh bất thường bốn kênh được chuyển tiếp từ Pha 1, $R_\psi$ biểu diễn hàm trích xuất đặc trưng của mạng ResNet cải tiến với tập tham số $\psi$, và $f_{defect}$ là vectơ đặc trưng hình thái đầu ra được sử dụng cho thuật toán phân cụm trực tuyến.

Vectơ đặc trưng $f_{defect}$ tập trung phản ánh các yếu tố cấu trúc bao gồm: (i) kích thước và diện tích của vùng khuyết tật, (ii) độ bất đồng nhất về mặt entropy tại khu vực biến dạng, và (iii) tỷ lệ hình dạng biên giới khuyết tật. Các yếu tố này tạo tiền đề cho việc phân tách chính xác các dạng lỗi có hình thái khác nhau trong không gian đặc trưng.

\subsection{Hàm mất mát và Tối ưu hóa mạng trích xuất đặc trưng}
\label{subsec:hammatmat_toiuu}

Quá trình huấn luyện mạng trích xuất đặc trưng sử dụng các hàm mất mát chuyên biệt nhằm tối ưu hóa khả năng biểu diễn dữ liệu cho từng pha. Việc lựa chọn đúng hàm mục tiêu đảm bảo mô hình không bị quá khớp với dữ liệu huấn luyện và duy trì tính ổn định khi triển khai thực tế.

Mạng mã hóa tự động ở Pha 1 được tối ưu hóa thông qua hàm mất mát sai số bình phương trung bình kết hợp với chỉ số độ tương đồng cấu trúc. Hàm mất mát tổng hợp của mạng mã hóa tự động được xác định theo Công thức \ref{eq:loss_ae}:
\begin{equation}
\label{eq:loss_ae}
L_{AE}(\theta, \phi) = \frac{1}{M} \sum_{i=1}^{M} \| I_{RGB-E}^{(i)} - \hat{I}_{RGB-E}^{(i)} \|_2^2 + \lambda \left( 1 - \text{SSIM}(I_{RGB-E}^{(i)}, \hat{I}_{RGB-E}^{(i)}) \right)
\end{equation}
trong đó $M$ đại diện cho tổng số mẫu ảnh trong một lô huấn luyện, $I_{RGB-E}^{(i)}$ và $\hat{I}_{RGB-E}^{(i)}$ lần lượt biểu diễn ảnh gốc thứ $i$ và ảnh tái tạo tương ứng, $\text{SSIM}$ là chỉ số đo lường độ tương đồng cấu trúc giữa hai hình ảnh, $\lambda$ là hệ số cân bằng giữa hai thành phần mất mát, và $\|\cdot\|_2$ thể hiện chuẩn Euclid hai.

Thuật toán tối ưu hóa Adam được áp dụng để cập nhật trọng số cho cả hai mạng với tốc độ học được điều chỉnh giảm dần theo thời gian huấn luyện. Quá trình tối ưu hóa kết thúc khi giá trị của hàm mất mát đạt trạng thái hội tụ và không giảm quá một ngưỡng xác định trước trong nhiều chu kỳ huấn luyện liên tiếp v.v…
\section{Mô hình phát hiện bất thường ngoại tuyến}
\label{sec:mo_hinh_bat_thuong_ngoai_tuyen}

Pha phát hiện bất thường ngoại tuyến đóng vai trò sàng lọc toàn bộ các sản phẩm màn hình điện thoại đạt tiêu chuẩn chất lượng trước khi đưa dữ liệu nghi ngờ vào chuỗi xử lý phức tạp ở pha tiếp theo. Việc xây dựng một mô hình đại diện chính xác cho phân bố của sản phẩm không khuyết tật giúp hệ thống tối ưu hóa tài nguyên tính toán và giảm thiểu tỷ lệ báo động giả. Cơ chế này dựa trên sự kết hợp giữa khả năng nén đặc trưng của mạng mã hóa tự động và năng lực học tô-pô không gian của Mạng khí Nơ-ron tăng trưởng.

\subsection{Huấn luyện GNG ngoại tuyến trên không gian tiềm ẩn}
\label{subsec:huan_luyen_gng_ngoai_tuyen}

Mạng khí nơ-ron tăng trưởng ngoại tuyến được xây dựng nhằm mô phỏng chính xác cấu trúc hình học của không gian đặc trưng biểu diễn bề mặt sản phẩm đạt chuẩn. Tập hợp các vectơ tiềm ẩn thu được từ đầu ra của mạng mã hóa tự động ở Mục \ref{subsec:autoencoder} đóng vai trò là dữ liệu đầu vào cho quá trình huấn luyện. Thông qua thuật toán cạnh tranh không giám sát, mạng tự động điều chỉnh vị trí các nút và thiết lập các liên kết cạnh nhằm bao phủ toàn bộ vùng dữ liệu chuẩn trong không gian tiềm ẩn \cite{fritzke1995growing}.

Quá trình cập nhật vị trí nút chiến thắng và các nút lân cận trực tiếp trong mạng khí nơ-ron tăng trưởng được thực hiện theo nguyên lý học cạnh tranh \cite{martinetz1993neural}. Đối với mỗi vectơ tiềm ẩn đầu vào $z$, nút chiến thắng thứ nhất $s_1$ có khoảng cách Euclid nhỏ nhất đến $z$ và nút chiến thắng thứ hai $s_2$ được xác định. Vectơ trọng số của nút $s_1$ được cập nhật theo Công thức \ref{eq:gng_update_winner}:
\begin{equation}
\label{eq:gng_update_winner}
\Delta w_{s_1} = \epsilon_b (z - w_{s_1})
\end{equation}
trong đó $z$ đại diện cho vectơ tiềm ẩn đầu vào, $w_{s_1}$ biểu diễn vectơ trọng số của nút chiến thắng thứ nhất $s_1$, và $\epsilon_b$ là tốc độ học của nút chiến thắng. Đồng thời, vectơ trọng số của các nút $n$ thuộc tập hợp lân cận trực tiếp $N_{s_1}$ của nút $s_1$ được điều chỉnh theo Công thức \ref{eq:gng_update_neighbor}:
\begin{equation}
\label{eq:gng_update_neighbor}
\Delta w_n = \epsilon_n (z - w_n) \quad \forall n \in N_{s_1}
\end{equation}
trong đó $w_n$ biểu diễn vectơ trọng số của nút lân cận $n$, $N_{s_1}$ là tập hợp các nút có liên kết cạnh trực tiếp với nút $s_1$, và $\epsilon_n$ đại diện cho tốc độ học của các nút lân cận.

Sự tích lũy lỗi cục bộ tại từng nút quyết định việc bổ sung nút mới vào cấu trúc mạng sau một số chu kỳ huấn luyện định kỳ. Lỗi tích lũy $A_{s_1}$ của nút chiến thắng $s_1$ được cập nhật theo Công thức \ref{eq:gng_error_accum}:
\begin{equation}
\label{eq:gng_error_accum}
A_{s_1} \leftarrow A_{s_1} + \| z - w_{s_1} \|_2^2
\end{equation}
trong đó $A_{s_1}$ đại diện cho giá trị lỗi tích lũy của nút $s_1$, $z$ là vectơ tiềm ẩn hiện tại, $w_{s_1}$ là vectơ trọng số của nút $s_1$, và $\| \cdot \|_2$ thể hiện chuẩn Euclid hai. Quá trình gia tăng số lượng nút dừng lại khi mạng đạt đến số lượng nút tối đa quy định trước hoặc khi sai số tái tạo toàn cục đạt trạng thái ổn định.

Mô hình phân bố dữ liệu chuẩn thu được sau quá trình huấn luyện ngoại tuyến hình thành nên một đồ thị phác họa vùng không gian an toàn của sản phẩm không lỗi. Sơ đồ minh họa quá trình phân tách vùng dữ liệu chuẩn và xác định điểm bất thường trên không gian tiềm ẩn được thể hiện trong Hình \ref{fig:gng_manifold}.

\begin{figure}[h!]
\centering
\begin{tikzpicture}[font=\normalsize, >=stealth]
    % 1. Miền không gian tiềm ẩn chuẩn Z_good
    \path[fill=blue!5, draw=blue!30, dashed, thick, rounded corners=20pt] 
        (-0.6,-0.6) rectangle (5.8,3.6);
    \node[black, font=\normalsize] at (2.6,4.3) {Miền không gian tiềm ẩn chuẩn ($Z_{\text{good}}$)};

    % 2. Tọa độ các nút GNG
    \coordinate (N1) at (0.5,0.6);
    \coordinate (N2) at (1.8,1.0); % s2
    \coordinate (N3) at (1.4,2.3); % s1 (Winner)
    \coordinate (N4) at (3.3,0.8);
    \coordinate (N5) at (3.6,2.3); % Neighbor of s1
    \coordinate (N6) at (4.8,1.4);

    % 3. Ranh giới quyết định (bao quanh các nút với bán kính mu + 3*sigma)
    \draw[red!70!black, dotted, line width=1pt, fill=red!5, fill opacity=0.3] 
        plot[smooth cycle, tension=0.6] coordinates {
            (-0.2,0.5) (0.8,2.7) (3.8,2.8) (5.3,1.5) (3.8,0.3) (0.8,0.2)
        };
    \node[black, font=\small, align=center] at (4.6,-0.3) {Ranh giới an toàn\\$\mu + 3\sigma$ ($T_{\text{pha1}}$)};

    % 4. Các cạnh nối GNG
    \draw[black!70, thick] (N1) -- (N2);
    \draw[black!70, thick] (N1) -- (N3);
    \draw[black!70, thick] (N2) -- (N3);
    \draw[black!70, thick] (N2) -- (N4);
    \draw[black!70, thick] (N3) -- (N5);
    \draw[black!70, thick] (N4) -- (N5);
    \draw[black!70, thick] (N4) -- (N6);
    \draw[black!70, thick] (N5) -- (N6);

    % 5. Vẽ các nút GNG thông thường
    \foreach \p in {N1,N4,N6}
        \fill[blue!70!black] (\p) circle (3pt);

    % Nút chiến thắng s1 và s2, lân cận N5
    \fill[orange!90!black] (N3) circle (3.5pt) node[above, xshift=3mm, font=\scriptsize, black] {$s_1$ (Winner)};
    \fill[black] (N2) circle (3pt) node[below, font=\scriptsize] {$s_2$};
    \fill[black] (N5) circle (3pt) node[above, font=\scriptsize] {$n \in N_{s_1}$};

    % 6. Mẫu dữ liệu bình thường z_normal
    \coordinate (Zn) at (2.2,1.8);
    \fill[black] (Zn) circle (2.5pt);
    \draw[->, thick, black] (Zn) -- (N3) 
        node[midway, above, xshift=7mm, yshift=-2mm, align=center] {$\scriptstyle d \le \mu+3\sigma$};
    \node[below, black, font=\small] at (2.3,1.7) {$z_{\text{normal}}$};

    % 7. Mẫu dữ liệu bất thường z_anomaly
    \coordinate (Za) at (7.8,2.5);
    \fill[black] (Za) circle (2.5pt);
    \draw[->, thick, black, dashed] (Za) -- (N5) 
        node[midway, above right, align=center, text=black] {$\scriptstyle d \le \mu+3\sigma$};
    \node[above right, black, font=\small] at (Za) {$z_{\text{anomaly}}$ (Bất thường)};

\end{tikzpicture}
\caption{Mô phỏng không gian tiềm ẩn chuẩn với cấu trúc đồ thị Mạng khí Nơ-ron tăng trưởng ngoại tuyến và cơ chế xác định mẫu bất thường dựa trên khoảng cách ngưỡng $\mu + 3\sigma$.}
\label{fig:gng_manifold}
\end{figure}

\subsection{Cơ chế phân loại bất thường dựa trên khoảng cách ngưỡng ba lần độ lệch chuẩn}
\label{subsec:phan_loai_bat_thuong_3sigma}

Cơ chế phân loại bất thường dựa trên nguyên lý đo lường sự độ lệch của biểu diễn không gian ẩn từ mẫu đầu vào so với cấu trúc đồ thị của mô hình Mạng khí Nơ-ron tăng trưởng ngoại tuyến. Khi hình ảnh bề mặt màn hình điện thoại được đưa vào hệ thống, véc-tơ biểu diễn ẩn $z_x$ của mẫu ảnh đó sẽ được trích xuất thông qua mạng mã hóa tự động. Khoảng cách Euclid nhỏ nhất từ véc-tơ ẩn $z_x$ tới tập hợp các véc-tơ trọng số trong đồ thị Mạng khí Nơ-ron tăng trưởng ngoại tuyến được tính toán theo Công thức \ref{eq:min_distance}:
\begin{equation}
\label{eq:min_distance}
d(z_x) = \min_{i \in V} \| z_x - w_i \|_2
\end{equation}
Trong Công thức \ref{eq:min_distance}, ký hiệu $d(z_x)$ đại diện cho khoảng cách nhỏ nhất từ mẫu kiểm thử $z_x$ đến đồ thị mạng, ký hiệu $V$ biểu diễn tập hợp chỉ số các nút thuộc mạng, ký hiệu $w_i$ thể hiện véc-tơ trọng số của nút thứ $i$, và ký hiệu $\| \cdot \|_2$ đại diện cho chuẩn Euclid cấp hai.

Cơ sở lý thuyết của việc định giá trị ngưỡng phân định bất thường xuất phát từ mối liên hệ giữa khoảng cách biểu diễn không gian và xác suất thuộc về phân phối của tập sản phẩm đạt chuẩn \cite{pukelsheim1994three, nist2012ehandbook}. Giá trị khoảng cách $d(z_x)$ từ véc-tơ đầu vào $z_x$ đến nút gần nhất càng lớn phản ảnh mẫu dữ liệu càng nằm xa vùng phân phối của dữ liệu đạt chất lượng đã được mô hình học trước đó. Giả định chuỗi khoảng cách cực tiểu của các mẫu đạt tiêu chuẩn tuân theo phân phối chuẩn, khoảng khoảng cách $\mu_d \pm 3\sigma_d$ bao phủ $99{,}73\%$ số lượng mẫu dữ liệu bình thường theo tiêu chuẩn thống kê \cite{nist2012ehandbook}. Do đó, bất kỳ mẫu sản phẩm nào có khoảng cách vượt quá ngưỡng trên sẽ được hệ thống gán nhãn là mẫu bất thường với độ tin cậy cao.

Để thiết lập ngưỡng phân định $T_{pha1}$ một cách tự động, hệ thống sử dụng tập kiểm định độc lập gồm các mẫu sản phẩm không chứa khuyết tật. Các thông số thống kê bao gồm giá trị trung bình và độ lệch chuẩn của chuỗi khoảng cách cực tiểu trên tập kiểm định được xác định lần lượt thông qua Công thức \ref{eq:mean_distance2} và Công thức \ref{eq:std_distance2}:
\begin{equation}
\label{eq:mean_distance2}
\mu_d = \frac{1}{K} \sum_{j=1}^{K} d(z_j^{val})
\end{equation}
\begin{equation}
\label{eq:std_distance2}
\sigma_d = \sqrt{\frac{1}{K-1} \sum_{j=1}^{K} \left( d(z_j^{val}) - \mu_d \right)^2}
\end{equation}
Trong Công thức \ref{eq:mean_distance2} và Công thức \ref{eq:std_distance2}, ký hiệu $K$ đại diện cho tổng số lượng mẫu kiểm định đạt chuẩn, ký hiệu $z_j^{val}$ biểu diễn véc-tơ ẩn của mẫu kiểm định thứ $j$, ký hiệu $\mu_d$ thể hiện giá trị trung bình mẫu của khoảng cách, và ký hiệu $\sigma_d$ đại diện cho độ lệch chuẩn mẫu.

Quy tắc xác định ngưỡng phân loại bất thường $T_{pha1}$ được thiết lập theo nguyên lý ba lần độ lệch chuẩn thể hiện trong Công thức \ref{eq:threshold_3sigma}:
\begin{equation}
\label{eq:threshold_3sigma}
T_{pha1} = \mu_d + 3 \sigma_d
\end{equation}
Trong Công thức \ref{eq:threshold_3sigma}, ký hiệu $T_{pha1}$ đóng vai trò là ngưỡng khoảng cách tối đa chấp nhận được đối với một sản phẩm đạt chuẩn.

Tiến trình ra quyết định kiểm tra chất lượng đối với một mẫu ảnh đầu vào được thực hiện phân định qua hai trường hợp cụ thể: (i) nếu khoảng cách $d(z_x) \le T_{pha1}$, mẫu dữ liệu được kết luận là sản phẩm đạt chuẩn và hệ thống kết thúc tiến trình xử lý đối với mẫu này, và (ii) nếu khoảng cách $d(z_x) > T_{pha1}$, mẫu dữ liệu bị xác định là bất thường và được chuyển tiếp sang Pha 2 để thực hiện phân tích đặc trưng và phân cụm trực tuyến. Cơ chế này đảm bảo chỉ những hình ảnh thực sự chứa khuyết tật mới đi vào tiến trình phân cụm trực tuyến, giúp tối ưu hóa chi phí tính toán và duy trì tốc độ vận hành cho toàn hệ thống.
\section{Thuật toán phân cụm lỗi trực tuyến EAGNG}
\label{sec:thuat_toan_eagng}

Thuật toán Mạng khí Nơ-ron tăng trưởng tích lũy bằng chứng được đề xuất nhằm giải quyết bài toán phân cụm khuyết tật thích nghi trực tuyến từ dòng đặc trưng trích xuất từ mạng ResNet bốn kênh ở Pha 2. Mạng khí Nơ-ron tăng trưởng truyền thống của Fritzke \cite{fritzke1995growing} gặp nhiều hạn chế khi vận hành trong môi trường công nghiệp liên tục do hiện tượng nhiễu dữ liệu, sự xuất hiện của các dạng lỗi mới và nguy cơ bùng nổ số lượng nút không kiểm soát. Thuật toán cải tiến tích hợp bốn cơ chế cốt lõi để duy trì tính ổn định tô-pô, kiểm soát mật độ nút và phân tách chính xác các hình thái khuyết tật phức tạp.

\subsection{Cơ chế suy giảm tốc độ học và tích lũy lỗi theo tuổi nút}
\label{subsec:node_aging_dynamic}

Cơ chế suy giảm tốc độ học và tích lũy lỗi theo tuổi nút điều chỉnh động năng lực học của từng nút dựa trên thời gian tồn tại của nút đó trong đồ thị. Các nút mới được tạo ra có mức độ linh hoạt cao để nhanh chóng thích nghi với vùng không gian đặc trưng mới, trong khi các nút có tuổi đời cao đạt trạng thái ổn định nhằm bảo tồn tri thức cấu trúc đã học \cite{martinetz1993neural}.

Tốc độ học của nút chiến thắng thứ nhất $s_1$ và các nút lân cận $n$ thuộc tập hợp $N_{s_1}$ giảm dần theo độ tuổi của nút $a_{s_1}$ được tính toán theo Công thức \ref{eq:aging_learning_rate_winner} và Công thức \ref{eq:aging_learning_rate_neighbor}:
\begin{equation}
\label{eq:aging_learning_rate_winner}
\epsilon_b(a_{s_1}) = \frac{\epsilon_{b0}}{1 + \alpha_a \cdot a_{s_1}}
\end{equation}
\begin{equation}
\label{eq:aging_learning_rate_neighbor}
\epsilon_n(a_n) = \frac{\epsilon_{n0}}{1 + \alpha_a \cdot a_n}
\end{equation}
trong đó $\epsilon_{b0}$ và $\epsilon_{n0}$ lần lượt đại diện cho tốc độ học khởi tạo của nút chiến thắng và nút lân cận, $a_{s_1}$ và $a_n$ tương ứng là tuổi hiện tại của nút $s_1$ và nút $n$, còn $\alpha_a$ là hệ số suy giảm tuổi nút.

Đồng thời, tốc độ tích lũy lỗi tại nút chiến thắng $s_1$ cũng được điều chỉnh giảm theo tuổi nhằm tránh hiện tượng phát triển nút mới bất hợp lý tại các vùng dữ liệu đã ổn định. Mức tăng lỗi tích lũy $\Delta A_{s_1}$ được xác định theo Công thức \ref{eq:error_aging_accum}:
\begin{equation}
\label{eq:error_aging_accum}
\Delta A_{s_1} = \left( 1 - \frac{a_{s_1}}{a_{s_1} + a_{max}} \right) \cdot \| f_{defect} - w_{s_1} \|_2^2
\end{equation}
trong đó $f_{defect}$ là vectơ đặc trưng lỗi đầu vào từ mạng ResNet, $w_{s_1}$ biểu diễn vectơ trọng số của nút $s_1$, $a_{max}$ là hằng số tuổi chuẩn hóa tối đa, và $\|\cdot\|_2$ thể hiện chuẩn Euclid hai. Cơ chế này đảm bảo những nút lâu năm không liên tục tích lũy lỗi do nhiễu ngẫu nhiên, giúp cấu trúc đồ thị duy trì sự ổn định qua thời gian.

\subsection{Kiểm tra chồng lấp Bounding Box kiểm soát liên kết cạnh}
\label{subsec:bounding_box_overlap}

Cơ chế kiểm tra chồng lấp hộp giới hạn ngăn chặn việc kết nối nhầm lẫn giữa hai cụm khuyết tật riêng biệt nằm gần nhau trong không gian đặc trưng. Phép tạo cạnh nối truyền thống chỉ dựa trên khoảng cách giữa hai nút chiến thắng có thể vô tình gộp hai loại lỗi có bản chất hình thái khác nhau thành một cụm duy nhất.

Đối với mỗi cụm nút $C_k$ trên đồ thị, một hộp giới hạn đa chiều $BB(C_k)$ được xác định bởi giá trị cực tiểu và cực đại của tọa độ các nút thuộc cụm đó. Tỷ lệ chồng lấp $O(C_a, C_b)$ giữa hai hộp giới hạn của cụm $C_a$ và cụm $C_b$ được tính toán theo Công thức \ref{eq:overlap_ratio2}:
\begin{equation}
\label{eq:overlap_ratio2}
O(C_a, C_b) = \frac{\text{Thể tích}(BB(C_a) \cap BB(C_b))}{\min \left( \text{Thể tích}(BB(C_a)), \text{Thể tích}(BB(C_b)) \right)}
\end{equation}
trong đó $\text{Thể tích}(\cdot)$ đại diện cho thể tích của hộp giới hạn trong không gian đặc trưng, và phần tử số biểu diễn thể tích phần giao nhau giữa hai hộp giới hạn.

Điều kiện thiết lập cạnh nối giữa nút $u \in C_a$ và nút $v \in C_b$ được quy định theo Công thức \ref{eq:edge_creation_rule}:
\begin{equation}
\label{eq:edge_creation_rule}
\text{Cho phép tạo cạnh }(u, v) \iff O(C_a, C_b) \ge \theta_{overlap}
\end{equation}
trong đó $\theta_{overlap}$ là ngưỡng chồng lấp tối thiểu bắt buộc. Nếu tỷ lệ chồng lấp $O(C_a, C_b)$ nhỏ hơn ngưỡng $\theta_{overlap}$, phép kết nối cạnh giữa hai nút bị loại bỏ, bảo toàn sự phân tách độc lập giữa hai lớp khuyết tật riêng biệt.


\subsection{Cơ chế trì hoãn chèn dựa trên bằng chứng và học cục bộ}
\label{subsec:deferred_insertion}

Cơ chế trì hoãn chèn dựa trên bằng chứng và học cục bộ giải quyết vấn đề chèn nút nhiễu khi dòng dữ liệu xuất hiện các điểm ngoại lệ đơn lẻ. Thay vì lập tức tạo nút mới khi gặp một mẫu dữ liệu nằm xa các nút hiện có, thuật toán lưu giữ mẫu dữ liệu đó vào một bộ đệm tạm thời để thu thập đủ bằng chứng mật độ trước khi tái cấu trúc tô-pô cục bộ \cite{penev1998local}.

Khi khoảng cách cực tiểu từ một vectơ đặc trưng mới $f_{defect}$ đến nút gần nhất vượt quá ngưỡng phân tách $T_{p2}$, điểm dữ liệu này được đưa vào vùng đệm lưu trữ tạm thời $B_r$ tương ứng với vùng không gian $r$. Ngưỡng quyết định việc khởi tạo quá trình huấn luyện cục bộ dựa trên số lượng bằng chứng tích lũy được xác định theo Công thức \ref{eq:buffer_threshold}:
\begin{equation}
\label{eq:buffer_threshold}
|B_r| \ge N_{evidence}
\end{equation}
trong đó $|B_r|$ đại diện cho số lượng mẫu đặc trưng hiện có trong bộ đệm $B_r$, và $N_{evidence}$ là hằng số ngưỡng bằng chứng tối thiểu bắt buộc để xác nhận một cụm lỗi mới.

Khi điều kiện trong Công thức \ref{eq:buffer_threshold} được thỏa mãn, thuật toán tiến hành quá trình huấn luyện cục bộ trên tập hợp các mẫu dữ liệu thuộc $B_r$. Trọng số của nút mới $w_{new}$ được khởi tạo tại tâm của tập bằng chứng $B_r$ theo Công thức \ref{eq:local_mini_train}:
\begin{equation}
\label{eq:local_mini_train}
w_{new} = \frac{1}{|B_r|} \sum_{x_j \in B_r} x_j
\end{equation}
trong đó $x_j$ đại diện cho vectơ đặc trưng thứ $j$ nằm trong bộ đệm $B_r$. Sau khi nút mới được khởi tạo và kết nối với các nút lân cận thích hợp, bộ đệm $B_r$ được xóa rỗng để sẵn sàng thu thập bằng chứng cho các chu kỳ tiếp theo v.v...

\subsection{Phát triển cấu trúc nhạy mật độ}
\label{subsec:density_aware_growth}

Cơ chế phát triển cấu trúc nhạy mật độ ngăn ngừa sự gia tăng số lượng nút quá mức tại các vùng đặc trưng có mật độ cao nhưng độ biến động nhỏ. Các thuật toán khí tăng trưởng truyền thống thường chèn nút mới tại vị trí có lỗi tích lũy lớn nhất mà không kiểm tra mật độ phân bố xung quanh, dẫn đến hiện tượng dư thừa nút và lãng phí bộ nhớ tính toán.

Mật độ nút cục bộ $\rho_q$ tại nút có lỗi tích lũy cao nhất $q$ được tính toán dựa trên khoảng cách trung bình đến các nút lân cận trực tiếp theo Công thức \ref{eq:local_density}:
\begin{equation}
\label{eq:local_density}
\rho_q = \frac{|N_q|}{\sum_{n \in N_q} \| w_q - w_n \|_2 + \epsilon_0}
\end{equation}
trong đó $N_q$ là tập hợp các nút lân cận trực tiếp có liên kết cạnh với nút $q$, $|N_q|$ đại diện cho số lượng nút lân cận, $w_q$ và $w_n$ lần lượt là vectơ trọng số của nút $q$ và nút $n$, còn $\epsilon_0$ là hằng số dương cực tiểu ngăn ngừa lỗi chia cho số không.

Quyết định chèn nút mới giữa nút $q$ và nút lân cận $f$ có lỗi lớn nhất được kiểm soát bởi điều kiện mật độ theo Công thức \ref{eq:density_check}:
\begin{equation}
\label{eq:density_check}
\text{Phép chèn được thực hiện khi } \rho_q < \rho_{thresh}
\end{equation}
trong đó $\rho_{thresh}$ là ngưỡng mật độ tối đa cho phép chèn nút. Nếu $\rho_q \ge \rho_{thresh}$, thuật toán bỏ qua việc chèn nút mới và thay vào đó thực hiện phân bổ lại lỗi tích lũy của nút $q$ cho các nút lân cận nhằm ổn định mạng mà không gia tăng quy mô đồ thị.


\subsection{Phân loại cụm lỗi bằng thuật toán Connected-Component trên đồ thị EAGNG}
\label{subsec:connected_component_clustering}

Sau khi cấu trúc đồ thị EAGNG được cập nhật trực tuyến, việc xác định các lớp khuyết tật dựa trên thuật toán tìm thành phần liên kết trên đồ thị vô hướng $G = (V, E)$. Mỗi thành phần liên kết cô lập đại diện cho một cụm khuyết tật có hình thái đặc trưng riêng biệt trên bề mặt màn hình điện thoại \cite{tarjan1972depth}.

Tập hợp các cụm khuyết tật $\mathcal{C} = \{C_1, C_2, \dots, C_K\}$ được trích xuất sao cho thỏa mãn hai điều kiện tô-pô cơ bản theo Công thức \ref{eq:cc_conditions}:
\begin{equation}
\label{eq:cc_conditions}
\bigcup_{i=1}^{K} C_i = V \quad \text{và} \quad C_i \cap C_j = \emptyset \quad \forall i \neq j
\end{equation}
trong đó $V$ đại diện cho tập hợp tất cả các nút trong đồ thị EAGNG, $E$ biểu diễn tập hợp các cạnh nối giữa các nút, và $K$ là tổng số cụm khuyết tật tự động phát hiện được tại thời điểm hiện tại.

Phương pháp phân cụm dựa trên đồ thị này loại bỏ yêu cầu phải khai báo trước số lượng cụm mục tiêu. Khi xuất hiện dạng khuyết tật hoàn toàn mới trên dây chuyền sản xuất, thuật toán EAGNG tự động hình thành các nút và thành phần liên kết mới, giúp hệ thống thích ứng linh hoạt với những biến đổi thực tế.

\subsection{Thuật toán tổng thể và Mã giả EAGNG}
\label{subsec:magia_eagng}

Thuật toán EAGNG tổng thể hợp nhất bốn cơ chế cải tiến vào một quy trình xử lý luồng dữ liệu liên tục. Quá trình vận hành của thuật toán được trình bày chi tiết trong Thuật toán \ref{alg:eagng_online}.

\bigskip
\noindent\hrule height 0.8pt
\vspace{2pt}
\captionsetup{justification=raggedright, singlelinecheck=false}
\captionof{algorithm}{Thuật toán phân cụm lỗi trực tuyến EAGNG}
\label{alg:eagng_online}
\vspace{-6pt}
\noindent\hrule
\vspace{4pt}

\begin{algorithmic}[1]

\Require Dòng đặc trưng $f_{defect}$, tập nút $V$, tập cạnh $E$, các tham số $\epsilon_{b0}, \epsilon_{n0}, N_{evidence}, \rho_{thresh}, \theta_{overlap}$.
\Ensure Tập các cụm khuyết tật $\mathcal{C} = \{C_1, C_2, \dots, C_K\}$.

\While{có mẫu dữ liệu đầu vào $f_{defect}$}

    \State Tìm hai nút $s_1, s_2 \in V$ gần $f_{defect}$ nhất theo chuẩn Euclid;

    \If{$\| f_{defect} - w_{s_1} \|_2 > T_{p2}$}

        \State Thêm $f_{defect}$ vào bộ đệm $B_r$;

        \If{$|B_r| \ge N_{evidence}$}
            \State Khởi tạo nút mới $w_{new}$ theo Công thức \ref{eq:local_mini_train};
            \State $V \leftarrow V \cup \{new\}$; Xóa bộ đệm $B_r$;
        \EndIf

    \Else

        \State Cập nhật tuổi tất cả các cạnh nối với $s_1$:
        $a_{(s_1,n)} \leftarrow a_{(s_1,n)} + 1$;

        \State Tính tốc độ học $\epsilon_b(a_{s_1})$ và $\epsilon_n(a_n)$
        theo Công thức \ref{eq:aging_learning_rate_winner}
        và \ref{eq:aging_learning_rate_neighbor};

        \State $w_{s_1} \leftarrow
        w_{s_1} + \epsilon_b(a_{s_1})
        (f_{defect} - w_{s_1})$;

        \State $w_n \leftarrow
        w_n + \epsilon_n(a_n)
        (f_{defect} - w_n),
        \quad \forall n \in N_{s_1}$;

        \State Tính tích lũy lỗi $\Delta A_{s_1}$
        theo Công thức \ref{eq:error_aging_accum};

        \If{chưa có cạnh $(s_1, s_2)$}

            \State Kiểm tra điều kiện chồng lấp
            $O(C_{s_1}, C_{s_2}) \ge \theta_{overlap}$
            theo Công thức \ref{eq:edge_creation_rule};

            \If{Thỏa mãn điều kiện}
                \State $E \leftarrow E \cup \{(s_1, s_2)\}$;
                \State Tái lập tuổi cạnh
                $a_{(s_1,s_2)} \leftarrow 0$;
            \EndIf

        \Else

            \State Tái lập tuổi cạnh
            $a_{(s_1,s_2)} \leftarrow 0$;

        \EndIf

        \If{đến chu kỳ kiểm tra chèn nút và $\rho_{s_1} < \rho_{thresh}$}
            \State Tìm nút lân cận $u \in N_{s_1}$ có lỗi lớn nhất
            và chèn nút mới giữa $s_1$ và $u$;
        \EndIf

        \State Loại bỏ các cạnh có tuổi vượt quá $a_{max}$
        và các nút cô lập không có liên kết;

    \EndIf

\EndWhile

\State Áp dụng thuật toán Connected-Component trích xuất tập cụm lỗi
$\mathcal{C}$ theo Công thức \ref{eq:cc_conditions};

\Return Tập cụm khuyết tật $\mathcal{C}$.

\end{algorithmic}
\vspace{2pt}
\noindent\hrule
\bigskip
\section{Công cụ tự động sinh tập dữ liệu khuyết tật mô phỏng}

Thuật toán phân cụm trực tuyến EAGNG yêu cầu một cơ chế kiểm thử tính thích ứng trước các biến đổi đột ngột và liên tục của dòng dữ liệu đầu vào. Do việc thu thập các hình ảnh lỗi thực tế trên dây chuyền sản xuất màn hình điện thoại thường gặp hạn chế về số lượng và sự đa dạng của các dạng khuyết tật, công cụ sinh dữ liệu tự động thích nghi đã được phát triển nhằm mục đích chủ động tạo ra các tập dữ liệu mô phỏng theo mô hình thống kê học \cite{gonzalez2018digital}. Mô hình này thực hiện việc tổng hợp ngẫu nhiên các cấu trúc nhiễu nền của màn hình chuẩn kết hợp với sự chèn lặp các dạng lỗi hình thái có kiểm soát về vị trí, kích thước và độ tương phản quang học.

Luận văn đề xuất quy trình tự động sinh dữ liệu hình ảnh được chia thành hai giai đoạn tính toán chính bao gồm giai đoạn mô phỏng nền màn hình và giai đoạn tổng hợp khuyết tật. Quy trình xử lý tổng thể của công cụ sinh dữ liệu được thể hiện trong Hình \ref{fig:generator_workflow}.

\begin{figure}[htbp]
\centering
\begin{tikzpicture}[node distance=1.0cm, auto, >=Latex, font=\normalsize]
  \tikzstyle{block} = [draw, fill=black!5, rectangle, minimum height=0.8cm, minimum width=3.2cm, text centered, rounded corners=2pt]
  \tikzstyle{line} = [draw, ->, thick]

  \node [block] (param) {Khởi tạo thông số hình học};
  \node [block, below=0.7cm of param] (bg) {Mô phỏng nhiễu nền phân phối chuẩn};
  \node [block, below=0.7cm of bg] (defect) {Tạo mặt nạ khuyết tật hình thái};
  \node [block, below=0.7cm of defect] (blur) {Lọc trung vị và hòa trộn quang học};
  \node [block, below=0.7cm of blur] (output) {Xuất ảnh kiểm thử màn hình};

  \path [line] (param) -- (bg);
  \path [line] (bg) -- (defect);
  \path [line] (defect) -- (blur);
  \path [line] (blur) -- (output);
\end{tikzpicture}
\caption{Mô hình quy trình tổng hợp hình ảnh khuyết tật màn hình mô phỏng.}
\label{fig:generator_workflow}
\end{figure}

Sơ đồ trong Hình \ref{fig:generator_workflow} minh họa các bước tính toán nối tiếp nhau nhằm tạo ra một hình ảnh màn hình hoàn chỉnh. Quá trình bắt đầu từ việc xác định các tham số không gian, tiếp tục qua việc tạo cấu trúc bề mặt bằng phân phối thống kê, sau đó vẽ mặt nạ khuyết tật và cuối cùng áp dụng bộ lọc mờ trung vị để đảm bảo tính liên tục của ranh giới điểm ảnh.

Trong giai đoạn mô phỏng nền màn hình, độ sáng của mỗi điểm ảnh tại tọa độ $(r, c)$ trên khung hình kích thước $H \times W$ điểm ảnh được tính toán kết hợp giữa thành phần biến ngẫu nhiên theo phân phối chuẩn và thành phần ngẫu nhiên động. Công thức (\ref{eq:bg_synthesis}) biểu diễn mô hình toán học tạo mức xám nền cho ảnh:

\begin{equation}
I_{bg}(r, c) = \text{clip}\left(0.8 \cdot \eta(r, c) + 0.2 \cdot \left[ B_{min} + U \cdot (B_{max} - B_{min}) \right], 0, 255\right)
\label{eq:bg_synthesis}
\end{equation}

Trong Công thức (\ref{eq:bg_synthesis}), ký hiệu $\eta(r, c) \sim \mathcal{N}(\mu, \sigma^2)$ đại diện cho mẫu ngẫu nhiên thu được từ phân phối chuẩn với giá trị trung bình $\mu = 215$ và độ lệch chuẩn $\sigma \in \{5, 7\}$, biến số $U \sim \mathcal{U}(0, 1)$ đại diện cho biến ngẫu nhiên phân phối đều trong khoảng từ zero đến một, và toán hàm $\text{clip}(x, 0, 255)$ thực hiện việc ép giá trị điểm ảnh nằm trong miền giới hạn mức xám từ 0 đến 255 theo chuẩn biểu diễn số nguyên tám bit. Các ngưỡng biên động $B_{min}$ và $B_{max}$ trong Công thức~(\ref{eq:bg_synthesis}) được xác định dựa trên Công thức~(\ref{eq:boundary_clamp}):

\begin{equation}
B_{min} = \max(0, \mu - 4\sigma), \quad B_{max} = \min(255, \mu + 3\sigma)
\label{eq:boundary_clamp}
\end{equation}

Công thức~(\ref{eq:boundary_clamp}) đảm bảo rằng thành phần độ sáng nền duy trì tính ổn định thống kê nhưng vẫn tạo ra sự biến thiên nhẹ giữa các khung hình nhằm mô phỏng nhiễu cảm biến quang học thực tế. Khi khởi tạo hình ảnh đạt tiêu chuẩn, giá trị độ lệch chuẩn $\sigma = 7$ được áp dụng để mô phỏng nhiễu phân tán rộng, trong khi hình ảnh chứa lỗi sử dụng giá trị $\sigma = 5$.

Đối với giai đoạn tạo khuyết tật, vị trí trung tâm $(x_d, y_d)$ của vùng lỗi được khống chế trong vùng không gian hiển thị an toàn nhằm tránh hiện tượng khuyết tật bị cắt đứt tại biên khung ảnh. Giới hạn tọa độ được mô tả qua Công thức~(\ref{eq:spatial_constraint}):

\begin{equation}
x_d \in \left[ \alpha_x W, (1 - \alpha_x) W \right], \quad y_d \in \left[ \alpha_y H, (1 - \alpha_y) H \right]
\label{eq:spatial_constraint}
\end{equation}

Trong Công thức~(\ref{eq:spatial_constraint}), thông số $\alpha_x = 0.15$ và $\alpha_y = 0.15$ đại diện cho tỷ lệ khoảng lùi tương đối so với chiều rộng $W = 540$ điểm ảnh và chiều cao $H = 1170$ điểm ảnh của màn hình. Kích thước bán kính đặc trưng của vết lỗi được tính toán dựa trên đường kính tối thiểu $D_{min} = \min(H, W)$ nhân với các tỷ lệ bán kính giới hạn trong khoảng từ $0.03$ đến $0.08$.

Công cụ hỗ trợ tổng hợp năm hình thái khuyết tật chính bao gồm: (i) hình tròn, (ii) hình tam giác, (iii) hình chữ nhật, (iv) đường thẳng ngang và (v) đường thẳng dọc. Mức độ nghiêm trọng của khuyết tật được phân thành bốn cấp độ thông qua hệ số tỷ lệ cường độ $\gamma_{level}$ bao gồm: (i) cấp độ rất thấp với $\gamma_{level} \in [0.01, 0.03]$, (ii) cấp độ thấp với $\gamma_{level} \in [0.03, 0.06]$, (iii) cấp độ trung bình với $\gamma_{level} \in [0.07, 0.11]$, và (iv) cấp độ cao với $\gamma_{level} \in [0.12, 0.20]$.

Sau khi vẽ mặt nạ khuyết tật binary $\mathbf{M}$, quá trình hòa trộn ảnh được thực hiện thông qua toán tử cộng có trọng số kết hợp với bộ lọc trung vị thích nghi. Công thức~(\ref{eq:defect_blend}) thể hiện phép toán tổng hợp hình ảnh cuối cùng:

\begin{equation}
I_{final}(r, c) = \text{clip}\left( I_{bg}(r, c) + (-1)^{\mathbb{I}_{dark}} \cdot \text{MedianBlur}\left( \mathbf{M}(r, c) \cdot \bar{I}_{bg} \cdot \gamma_{level}, k_{size} \right), 0, 255 \right)
\label{eq:defect_blend}
\end{equation}

Trong Công thức~(\ref{eq:defect_blend}), ký hiệu $\bar{I}_{bg}$ biểu thị giá trị độ sáng trung bình của ảnh nền, biến chỉ thị $\mathbb{I}_{dark} \in \{0, 1\}$ xác định tính chất lỗi tối hay lỗi sáng, và $\text{MedianBlur}(\cdot, k_{size})$ đại diện cho phép lọc mờ trung vị với kích thước cửa sổ $k_{size} = \max(3, 2 \lfloor 0.005 D_{min} \rfloor + 1)$ nhằm triệt tiêu các góc cạnh sắc nhọn nhân tạo, giúp hình thái vết lỗi tiệm cận với khuyết tật vật lý tự nhiên \cite{opencv2020manual}.

Mã nguồn thực thi quá trình sinh dữ liệu ảnh thích nghi được cài đặt với thư viện xử lý ảnh OpenCV được thể hiện chi tiết trong Thuật toán~\ref{alg:adaptive_generator}.

\bigskip
\noindent\hrule height 0.8pt
\vspace{2pt}
\captionsetup{justification=raggedright, singlelinecheck=false}
\captionof{algorithm}{Thuật toán sinh ảnh khuyết tật thích nghi.}
\label{alg:adaptive_generator}
\vspace{-6pt}
\noindent\hrule
\vspace{4pt}

\begin{algorithmic}[1]

\Require
Ảnh đầu vào $I$, phân phối chuẩn $\mathcal{N}(\mu,\sigma^2)$,
loại khuyết tật $T$, mức độ $L$, trạng thái lỗi sáng $B$

\Ensure
Ảnh đầu ra $I_{\mathrm{final}}$

\State Khởi tạo ma trận điểm ảnh $P$ kích thước $H \times W$

\State
$B_{\min} \gets \max(0,\mu-4\sigma)$

\State
$B_{\max} \gets \min(255,\mu+3\sigma)$

\For{$r \gets 0$ \textbf{to} $H-1$}
    \For{$c \gets 0$ \textbf{to} $W-1$}
        \State
        $s \gets \mathrm{NormalRandom}(\mu,\sigma)$
        
        \State
        $u \gets \mathrm{UniformRandom}(0,1)$
        
        \State
        $d \gets B_{\min}+u(B_{\max}-B_{\min})$
        
        \State
        $P[r,c] \gets
        \mathrm{Clip}(0.8s+0.2d,0,255)$
    \EndFor
\EndFor

\State
$I_{\mathrm{bg}} \gets \mathrm{CreateImage}(P)$

\If{$I$ chứa khuyết tật}
    \State
    $(x_d,y_d) \gets \mathrm{CalculateBoundedPoint}()$
    
    \State
    $\rho \gets \mathrm{GetScaleRatio}(L)$
    
    \State
    $A \gets \mathrm{Mean}(I_{\mathrm{bg}})\times\rho$
    
    \State
    $M \gets \mathrm{DrawDefectMask}(T,x_d,y_d,A)$
    
    \State
    $k \gets \mathrm{GetAdaptiveKernelSize}()$
    
    \State
    $M \gets \mathrm{MedianBlur}(M,k)$
    
    \If{$B = \mathrm{TRUE}$}
        \State $w \gets 1$
    \Else
        \State $w \gets -1$
    \EndIf
    
    \State
    $I_{\mathrm{final}}
    \gets \mathrm{Clip}(I_{\mathrm{bg}}+wM,0,255)$
\Else
    \State
    $I_{\mathrm{final}} \gets I_{\mathrm{bg}}$
\EndIf

\State \Return $I_{\mathrm{final}}$

\end{algorithmic}
\vspace{2pt}
\noindent\hrule
\bigskip

Thuật toán~\ref{alg:adaptive_generator} thể hiện sự tương quan trực tiếp với các mô hình toán học đã được thiết lập. Quá trình quét điểm ảnh hai chiều tạo ra cấu trúc nền tuân theo Công thức~(\ref{eq:bg_synthesis}), trong khi các hàm tính toán trọng số cường độ và lọc mờ trung vị thực hiện chính xác các phép toán được mô tả trong Công thức~(\ref{eq:defect_blend}). Việc tích hợp dịch vụ sinh ảnh tự động này cho phép hệ thống chủ động tạo ra hàng chục nghìn mẫu kiểm thử với các kịch bản trôi dạt khái niệm phức tạp, tạo tiền đề vững chắc cho việc đánh giá độ bền vững của thuật toán EAGNG trong các chương tiếp theo.



\section*{Kết luận Chương 3}
\addcontentsline{toc}{section}{Kết luận Chương 3}

Chương \ref{chap:c3} đã xây dựng hoàn chỉnh cơ sở lý thuyết, kiến trúc hệ thống và thuật toán đề xuất cho bài toán phát hiện bất thường và phân cụm khuyết tật màn hình điện thoại. Bằng việc kết hợp hai pha xử lý nối tiếp, hệ thống đã tối ưu hóa được năng lực tính toán thông qua cơ chế lọc bỏ mẫu đạt chuẩn ở Pha 1 trước khi tiến hành phân loại chuyên sâu ở Pha 2. Kỹ thuật ghép kênh entropy cục bộ vào dữ liệu ảnh nguyên bản đã nâng cao độ nhạy của hệ thống đối với các khuyết tật có độ tương phản thấp trên bề mặt sản phẩm.

Đóng góp lý thuyết quan trọng nhất của chương nằm ở thuật toán Mạng khí Nơ-ron tăng trưởng tích lũy bằng chứng với bốn cơ chế cải tiến đồng bộ. Sự kết hợp giữa động lực tuổi nút, cơ chế hoãn chèn dựa trên bằng chứng, kiểm soát mật độ phát triển nút và điều kiện chồng lấp hộp giới hạn đã giải quyết triệt để các hạn chế về nhiễu dữ liệu và bùng nổ cấu trúc của mô hình truyền thống. Kết quả phân cụm dựa trên thành phần liên kết đồ thị cho phép hệ thống tự động thích ứng với các dạng lỗi mới phát sinh mà không cần quy định trước số lượng cụm khuyết tật. Các mô hình toán học và thiết kế thuật toán trong chương này là nền tảng để tiến hành xây dựng kịch bản thực nghiệm, đánh giá định lượng và đối sánh hiệu năng trong Chương \ref{chap:c4}.